%---------------------------------------------------------------------------%
%->> 封面信息及生成
%---------------------------------------------------------------------------%
%-
%-> 中文封面信息
%-
\confidential{}% 密级：只有涉密论文才填写
\schoollogo{scale=0.095}{ucas_logo}% 校徽
\title{面向网络流量的重复音视频文件实时检测系统研究}% 论文中文题目
\author{陈冠林}% 论文作者
\advisor{刘庆云~正高级工程师}% 指导教师：姓名 专业技术职务 工作单位
\advisors{中国科学院信息工程研究所}% 指导老师附加信息 或 第二指导老师信息
\degree{工程硕士}% 学位：学士、硕士、博士
\major{计算机技术}% 二级学科专业名称
\institute{中国科学院信息工程研究所}% 院系名称
\date{2019~年~6~月}% 毕业日期：夏季为6月、冬季为12月
%-
%-> 英文封面信息
%-
\TITLE{Research on Real-time Detection System \\ of \\ Repeated Multimedia Files for Network Traffic}% 论文英文题目
\AUTHOR{Chen Guanlin}% 论文作者
\ADVISOR{Supervisor: Professor Liu Qingyun}% 指导教师
\DEGREE{Master}% 学位：Bachelor, Master, Doctor。封面格式将根据英文学位名称自动切换，请确保拼写准确无误
\DEGREETYPE{Engineering}% 学位类别：Philosophy, Natural Science, Engineering, Economics, Agriculture 等
\MAJOR{Computer Technology}% 二级学科专业名称
\INSTITUTE{Institute of Information Engineering, Chinese Academy of Sciences}% 院系名称
\DATE{June 2019}% 毕业日期：夏季为June、冬季为December
%-
%-> 生成封面
%-
\maketitle% 生成中文封面
\MAKETITLE% 生成英文封面
%-
%-> 作者声明
%-
\makedeclaration% 生成声明页
%-
%-> 中文摘要
%-
\chapter*{摘\quad 要}\chaptermark{摘\quad 要}% 摘要标题
\setcounter{page}{1}% 开始页码
\pagenumbering{Roman}% 页码符号

随着互联网的发展，网络已经渗入到我们生活的方方面面，越来越多的人使用网络观看和下载音视频节目，这也使得音视频流量已经成为网络流量中占比最高的部分。根据Sandvine发布的2018年全球互联网观察显示，视频流量已经占据全球流量的57.69\%，较上一年增长了22.43\%。

据统计，YouTube流量中约有30\%的重复。大量重复的音视频文件通过网络进行传输，而代理服务器、入侵检测系统等网络中间件需要对这些流量处理，以达到版权保护、色情检测等目的。网络中间件对于重复流量的处理消耗了大量资源。

为了解决网络中间件对重复音视频文件处理资源消耗的问题，本文构建了一套实时的重复音视频文件检测系统，该系统可以对网络流量中的重复音视频内容进行检测，用以指导网络中间件的处理，最终达到节约网络、存储和计算资源的目的。

该系统使用以下技术：1）预测性文件标识生成方法。该技术采用信息熵和互信息作为评价标准，对生成文件标识的特征进行筛选，并在流量建立连接时，生成一个具有预测性的文件标识。 2）基于流式模糊哈希的重复音视频检测方法。该技术通过预测性文件标识检测文件的重复性，同时计算流式模糊哈希值生成文件编号，以验证标识的效果，并指导标识的特征选择。

本文的创新点包括以下两点：

1）提出了一种预测性文件标识生成方法，该方法在音视频传输初始阶段生成音视频文件标识。

2）提出了一种基于流式模糊哈希的重复音视频检测方法，该方法可以检测完整的音视频文件的重复性，并指导预测性文件标识的生成。

\keywords{网络流，重复文件，识别方法，预测性}% 中文关键词
%-
%-> 英文摘要
%-
\chapter*{Abstract}\chaptermark{Abstract}% 摘要标题

With the development of the Internet, the network has penetrated into every aspect of our lives. More and more people use the network to watch and download multimedia program, which makes multimedia traffic account for the highest proportion of network traffic. According to the 2018 internet phenomena report released by Sandvine, video traffic has accounted for 57.69\% of global traffic and increased 22.43\% compared with last year.

According to the statistics, around 30\% of the YouTube traffic is identified as redundant. There is a large number of repeated files in these multimedia traffic, and network middlebox like proxy servers and intrusion detection systems consume a lot of resources for processing these repeated multimedia traffic. In this paper, we construct a real-time repeated multimedia file detection system, which can detect the repeated multimedia files in network traffic and guide the processing of network middlebox.

This system uses the following two technologies: 1) Predictive file identification generation method. This technology uses information entropy and mutual information as evaluation criteria filter the features of the generated file identifiers and generates a predictive file identification when the traffic establishes a connection. 2) A method of detecting repeated multimedia based on streaming fuzzy hashing. If the file identification is queried in the database, the multimedia file is considered to be repeated, and the file number is generated by the stream fuzzy hash value to verify the effect of the file identification, and the result can be used to guide the feature selection of the file identification. 

The innovations of this paper include the following two points:

1) A predictive file identification generation method is proposed, which detects repeated multimedia files in the initial stage of multimedia transmission.

2) A method of detecting the repeated multimedia files based on streaming fuzzy hash is proposed. This method can detect the repeated multimedia files, and the detection results can be used to guide feature selection of the predictive file identification generation.

\KEYWORDS{Network Traffic, Repeated File, Identification Method, Predictive}% 英文关键词
%---------------------------------------------------------------------------%
